\section{System Architecture}
\label{sec:arch}
\system{} integrates multi-agent reasoning with industrial supervision around a common representation of production assets and operating context. Its organizing object is the node: an acquisition or control endpoint with engineering semantics, device connections, line membership, and agent bindings. A channel organizes the agents that interpret these endpoints, controllers mediate plant interactions, and Fig.~\ref{fig:arch} locates these components and execution boundaries.

\begin{figure*}[!tb]
\centering
\includegraphics[width=0.76\textwidth]{figures/publication/fig2-architecture.pdf}
\caption{System architecture and execution boundaries. The co-resident platform core connects agent channels, semantic nodes, industrial controllers, and operator services. CLI/OMP engines and storage services extend beyond that process. Shared identifiers support attribution across the logical data model; they do not make storage or physical actuation transactional.}
\label{fig:arch}
\end{figure*}

\subsection{A Shared Supervisory Core}
The Node.js core hosts channel routing, task scheduling, industrial services,
protocol adapters, and platform interfaces: REST and WebSocket endpoints serve
operator applications, MCP and A2A interfaces support external integration, and a
plugin host extends platform capabilities. These surfaces reuse platform objects,
but their exposed operations and permission checks differ, and an agent's industrial
authority derives from its node bindings rather than from channel access or the
availability of a protocol endpoint. Co-residency reduces the need to synchronize
duplicate nodes, bindings, and runs across separately assembled applications, but
does not place the deployment in one process: CLI engines and the platform's bundled
OMP agent harness execute in spawned processes, TimescaleDB and MinIO are external,
and connection establishment, queued protocol operations, and readback introduce
asynchronous boundaries. Fast regulation and protection therefore remain with the
plant controller and its safety equipment.

\subsection{Reasoning, Plant Interaction, and Operational Context}
On the reasoning side, a channel combines a lead, worker instances, a task tree,
messages, and retrievable memory. Engine adapters allow different reasoning
implementations to participate without changing a bound node's identity or the
controller interface, which separates task organization from device authority.

On the plant side, acquisition nodes turn protocol readings into engineering
quantities, control nodes translate engineering setpoints into device commands and
readback results, per-node runtimes manage timing and in-flight operations, and
controller services supply gateway state, recipe context, and write bookkeeping.
Nodes and devices have many-to-many configuration associations, but association is
distinct from execution routing: a control-node write invokes its configured driver
and connection settings rather than iterating over its associated devices, so the
interface follows process meaning without implying multi-device actuation by one
write.

The bridge between these sides consists of semantic cards, instance-specific
bindings, industrial tools, and source-dependent controller checks. Bindings
determine whether a channel member can request an operation and whether human
approval is needed; controllers then apply the checks appropriate to that request's
origin, so semantic guidance helps an agent propose a sensible action whereas
executable checks determine which requests reach a driver. Prompt guidance, such as
the suggested step, is not an enforced interlock.

\subsection{Logical Continuity Across a Production Run}
Lines, products, recipes, and runs provide the context in which measurements and
interventions become comparable. Their identifiers connect platform state, time
series, image frames, operational events, and optimization records across
heterogeneous stores: SQLite and file-backed repositories for platform state,
TimescaleDB for measurements, MinIO for objects. These constitute one logical data
model rather than one transaction domain. An interaction begins when an engineer
connects devices, defines nodes, associates production context, and binds channel
members; starting a line establishes an active run and initiates its parameter
writes. Agents then inspect observations, exchange evidence, and propose adjustments
through their bound tools; accepted requests proceed to device communication and
readback, and subsequent observations inform an agent-supplied judgment and an
independent rule-based backstop, either of which may require a further physical
write.

Operators participate through approvals, line operations, and monitoring.
Section~\ref{sec:mech} details these mechanisms and marks where their guarantees
end.
